\documentclass[12pt]{article}
\usepackage[utf8]{inputenc}
\usepackage{listings}
\usepackage{xcolor}
\usepackage{geometry}
\usepackage{parskip}

\geometry{margin=2.5cm}

\definecolor{codebg}{RGB}{245,245,245}
\definecolor{keyword}{RGB}{0,112,0}
\definecolor{comment}{RGB}{128,128,128}
\definecolor{string}{RGB}{180,0,0}
\definecolor{outputbg}{RGB}{240,248,255}

\lstdefinestyle{pythonstyle}{
    backgroundcolor=\color{codebg},
    basicstyle=\ttfamily\small,
    keywordstyle=\color{keyword}\bfseries,
    stringstyle=\color{string},
    commentstyle=\color{comment}\itshape,
    language=Python,
    frame=single,
    rulecolor=\color{gray!40},
    breaklines=true,
    showstringspaces=false,
    tabsize=4,
    captionpos=b,
}

\lstdefinestyle{outputstyle}{
    backgroundcolor=\color{outputbg},
    basicstyle=\ttfamily\small,
    frame=single,
    rulecolor=\color{blue!20},
    breaklines=true,
}

\title{\textbf{Wind Turbine Blade Fault Classification\\Using a Simple CNN}\\[0.5em]
\large Step 3 --- Print Dataset Information}
\author{Amreen Batool}
\date{}

\begin{document}
\maketitle

\section*{Step 3 --- Print Dataset Information}

\begin{lstlisting}[style=pythonstyle]
def count_images(folder):
    total = 0
    class_counts = {}

    for class_name in sorted(os.listdir(folder)):
        class_path = os.path.join(folder, class_name)
        if os.path.isdir(class_path):
            num_images = len([
                f for f in os.listdir(class_path)
                if f.lower().endswith(('.png', '.jpg', '.jpeg'))
            ])
            class_counts[class_name] = num_images
            total += num_images

    return total, class_counts

train_total, train_counts = count_images(train_dir)
val_total,   val_counts   = count_images(val_dir)
test_total,  test_counts  = count_images(test_dir)

print("Training set total:",   train_total)
print("Training class counts:", train_counts)
print()
print("Validation set total:",   val_total)
print("Validation class counts:", val_counts)
print()
print("Test set total:",   test_total)
print("Test class counts:", test_counts)
\end{lstlisting}

\subsection*{Output}

\begin{lstlisting}[style=outputstyle]
Training set total: 4200
Training class counts: {'Faulty': 2100, 'Healthy': 2100}

Validation set total: 600
Validation class counts: {'Faulty': 300, 'Healthy': 300}

Test set total: 1200
Test class counts: {'Faulty': 600, 'Healthy': 600}
\end{lstlisting}

\subsection*{Explanation}

This block counts:
\begin{itemize}
    \item Total number of images in each split (Train / Validation / Test).
    \item Number of \texttt{Faulty} and \texttt{Healthy} images in each split.
\end{itemize}

\subsection*{Why It Is Useful}

You can use this output in your report to describe the dataset composition and confirm that the classes are balanced.

\end{document}
